% Options for packages loaded elsewhere
% Options for packages loaded elsewhere
\PassOptionsToPackage{unicode}{hyperref}
\PassOptionsToPackage{hyphens}{url}
\PassOptionsToPackage{dvipsnames,svgnames,x11names}{xcolor}
%
\documentclass[
  letterpaper,
  DIV=11,
  numbers=noendperiod]{scrreprt}
\usepackage{xcolor}
\usepackage{amsmath,amssymb}
\setcounter{secnumdepth}{5}
\usepackage{iftex}
\ifPDFTeX
  \usepackage[T1]{fontenc}
  \usepackage[utf8]{inputenc}
  \usepackage{textcomp} % provide euro and other symbols
\else % if luatex or xetex
  \usepackage{unicode-math} % this also loads fontspec
  \defaultfontfeatures{Scale=MatchLowercase}
  \defaultfontfeatures[\rmfamily]{Ligatures=TeX,Scale=1}
\fi
\usepackage{lmodern}
\ifPDFTeX\else
  % xetex/luatex font selection
\fi
% Use upquote if available, for straight quotes in verbatim environments
\IfFileExists{upquote.sty}{\usepackage{upquote}}{}
\IfFileExists{microtype.sty}{% use microtype if available
  \usepackage[]{microtype}
  \UseMicrotypeSet[protrusion]{basicmath} % disable protrusion for tt fonts
}{}
\makeatletter
\@ifundefined{KOMAClassName}{% if non-KOMA class
  \IfFileExists{parskip.sty}{%
    \usepackage{parskip}
  }{% else
    \setlength{\parindent}{0pt}
    \setlength{\parskip}{6pt plus 2pt minus 1pt}}
}{% if KOMA class
  \KOMAoptions{parskip=half}}
\makeatother
% Make \paragraph and \subparagraph free-standing
\makeatletter
\ifx\paragraph\undefined\else
  \let\oldparagraph\paragraph
  \renewcommand{\paragraph}{
    \@ifstar
      \xxxParagraphStar
      \xxxParagraphNoStar
  }
  \newcommand{\xxxParagraphStar}[1]{\oldparagraph*{#1}\mbox{}}
  \newcommand{\xxxParagraphNoStar}[1]{\oldparagraph{#1}\mbox{}}
\fi
\ifx\subparagraph\undefined\else
  \let\oldsubparagraph\subparagraph
  \renewcommand{\subparagraph}{
    \@ifstar
      \xxxSubParagraphStar
      \xxxSubParagraphNoStar
  }
  \newcommand{\xxxSubParagraphStar}[1]{\oldsubparagraph*{#1}\mbox{}}
  \newcommand{\xxxSubParagraphNoStar}[1]{\oldsubparagraph{#1}\mbox{}}
\fi
\makeatother

\usepackage{color}
\usepackage{fancyvrb}
\newcommand{\VerbBar}{|}
\newcommand{\VERB}{\Verb[commandchars=\\\{\}]}
\DefineVerbatimEnvironment{Highlighting}{Verbatim}{commandchars=\\\{\}}
% Add ',fontsize=\small' for more characters per line
\usepackage{framed}
\definecolor{shadecolor}{RGB}{241,243,245}
\newenvironment{Shaded}{\begin{snugshade}}{\end{snugshade}}
\newcommand{\AlertTok}[1]{\textcolor[rgb]{0.68,0.00,0.00}{#1}}
\newcommand{\AnnotationTok}[1]{\textcolor[rgb]{0.37,0.37,0.37}{#1}}
\newcommand{\AttributeTok}[1]{\textcolor[rgb]{0.40,0.45,0.13}{#1}}
\newcommand{\BaseNTok}[1]{\textcolor[rgb]{0.68,0.00,0.00}{#1}}
\newcommand{\BuiltInTok}[1]{\textcolor[rgb]{0.00,0.23,0.31}{#1}}
\newcommand{\CharTok}[1]{\textcolor[rgb]{0.13,0.47,0.30}{#1}}
\newcommand{\CommentTok}[1]{\textcolor[rgb]{0.37,0.37,0.37}{#1}}
\newcommand{\CommentVarTok}[1]{\textcolor[rgb]{0.37,0.37,0.37}{\textit{#1}}}
\newcommand{\ConstantTok}[1]{\textcolor[rgb]{0.56,0.35,0.01}{#1}}
\newcommand{\ControlFlowTok}[1]{\textcolor[rgb]{0.00,0.23,0.31}{\textbf{#1}}}
\newcommand{\DataTypeTok}[1]{\textcolor[rgb]{0.68,0.00,0.00}{#1}}
\newcommand{\DecValTok}[1]{\textcolor[rgb]{0.68,0.00,0.00}{#1}}
\newcommand{\DocumentationTok}[1]{\textcolor[rgb]{0.37,0.37,0.37}{\textit{#1}}}
\newcommand{\ErrorTok}[1]{\textcolor[rgb]{0.68,0.00,0.00}{#1}}
\newcommand{\ExtensionTok}[1]{\textcolor[rgb]{0.00,0.23,0.31}{#1}}
\newcommand{\FloatTok}[1]{\textcolor[rgb]{0.68,0.00,0.00}{#1}}
\newcommand{\FunctionTok}[1]{\textcolor[rgb]{0.28,0.35,0.67}{#1}}
\newcommand{\ImportTok}[1]{\textcolor[rgb]{0.00,0.46,0.62}{#1}}
\newcommand{\InformationTok}[1]{\textcolor[rgb]{0.37,0.37,0.37}{#1}}
\newcommand{\KeywordTok}[1]{\textcolor[rgb]{0.00,0.23,0.31}{\textbf{#1}}}
\newcommand{\NormalTok}[1]{\textcolor[rgb]{0.00,0.23,0.31}{#1}}
\newcommand{\OperatorTok}[1]{\textcolor[rgb]{0.37,0.37,0.37}{#1}}
\newcommand{\OtherTok}[1]{\textcolor[rgb]{0.00,0.23,0.31}{#1}}
\newcommand{\PreprocessorTok}[1]{\textcolor[rgb]{0.68,0.00,0.00}{#1}}
\newcommand{\RegionMarkerTok}[1]{\textcolor[rgb]{0.00,0.23,0.31}{#1}}
\newcommand{\SpecialCharTok}[1]{\textcolor[rgb]{0.37,0.37,0.37}{#1}}
\newcommand{\SpecialStringTok}[1]{\textcolor[rgb]{0.13,0.47,0.30}{#1}}
\newcommand{\StringTok}[1]{\textcolor[rgb]{0.13,0.47,0.30}{#1}}
\newcommand{\VariableTok}[1]{\textcolor[rgb]{0.07,0.07,0.07}{#1}}
\newcommand{\VerbatimStringTok}[1]{\textcolor[rgb]{0.13,0.47,0.30}{#1}}
\newcommand{\WarningTok}[1]{\textcolor[rgb]{0.37,0.37,0.37}{\textit{#1}}}

\usepackage{longtable,booktabs,array}
\newcounter{none} % for unnumbered tables
\usepackage{calc} % for calculating minipage widths
% Correct order of tables after \paragraph or \subparagraph
\usepackage{etoolbox}
\makeatletter
\patchcmd\longtable{\par}{\if@noskipsec\mbox{}\fi\par}{}{}
\makeatother
% Allow footnotes in longtable head/foot
\IfFileExists{footnotehyper.sty}{\usepackage{footnotehyper}}{\usepackage{footnote}}
\makesavenoteenv{longtable}
\usepackage{graphicx}
\makeatletter
\newsavebox\pandoc@box
\newcommand*\pandocbounded[1]{% scales image to fit in text height/width
  \sbox\pandoc@box{#1}%
  \Gscale@div\@tempa{\textheight}{\dimexpr\ht\pandoc@box+\dp\pandoc@box\relax}%
  \Gscale@div\@tempb{\linewidth}{\wd\pandoc@box}%
  \ifdim\@tempb\p@<\@tempa\p@\let\@tempa\@tempb\fi% select the smaller of both
  \ifdim\@tempa\p@<\p@\scalebox{\@tempa}{\usebox\pandoc@box}%
  \else\usebox{\pandoc@box}%
  \fi%
}
% Set default figure placement to htbp
\def\fps@figure{htbp}
\makeatother


% definitions for citeproc citations
\NewDocumentCommand\citeproctext{}{}
\NewDocumentCommand\citeproc{mm}{%
  \begingroup\def\citeproctext{#2}\cite{#1}\endgroup}
\makeatletter
 % allow citations to break across lines
 \let\@cite@ofmt\@firstofone
 % avoid brackets around text for \cite:
 \def\@biblabel#1{}
 \def\@cite#1#2{{#1\if@tempswa , #2\fi}}
\makeatother
\newlength{\cslhangindent}
\setlength{\cslhangindent}{1.5em}
\newlength{\csllabelwidth}
\setlength{\csllabelwidth}{3em}
\newenvironment{CSLReferences}[2] % #1 hanging-indent, #2 entry-spacing
 {\begin{list}{}{%
  \setlength{\itemindent}{0pt}
  \setlength{\leftmargin}{0pt}
  \setlength{\parsep}{0pt}
  % turn on hanging indent if param 1 is 1
  \ifodd #1
   \setlength{\leftmargin}{\cslhangindent}
   \setlength{\itemindent}{-1\cslhangindent}
  \fi
  % set entry spacing
  \setlength{\itemsep}{#2\baselineskip}}}
 {\end{list}}
\usepackage{calc}
\newcommand{\CSLBlock}[1]{\hfill\break\parbox[t]{\linewidth}{\strut\ignorespaces#1\strut}}
\newcommand{\CSLLeftMargin}[1]{\parbox[t]{\csllabelwidth}{\strut#1\strut}}
\newcommand{\CSLRightInline}[1]{\parbox[t]{\linewidth - \csllabelwidth}{\strut#1\strut}}
\newcommand{\CSLIndent}[1]{\hspace{\cslhangindent}#1}



\setlength{\emergencystretch}{3em} % prevent overfull lines

\providecommand{\tightlist}{%
  \setlength{\itemsep}{0pt}\setlength{\parskip}{0pt}}



 


\KOMAoption{captions}{tableheading}
\makeatletter
\@ifpackageloaded{tcolorbox}{}{\usepackage[skins,breakable]{tcolorbox}}
\@ifpackageloaded{fontawesome5}{}{\usepackage{fontawesome5}}
\definecolor{quarto-callout-color}{HTML}{909090}
\definecolor{quarto-callout-note-color}{HTML}{0758E5}
\definecolor{quarto-callout-important-color}{HTML}{CC1914}
\definecolor{quarto-callout-warning-color}{HTML}{EB9113}
\definecolor{quarto-callout-tip-color}{HTML}{00A047}
\definecolor{quarto-callout-caution-color}{HTML}{FC5300}
\definecolor{quarto-callout-color-frame}{HTML}{acacac}
\definecolor{quarto-callout-note-color-frame}{HTML}{4582ec}
\definecolor{quarto-callout-important-color-frame}{HTML}{d9534f}
\definecolor{quarto-callout-warning-color-frame}{HTML}{f0ad4e}
\definecolor{quarto-callout-tip-color-frame}{HTML}{02b875}
\definecolor{quarto-callout-caution-color-frame}{HTML}{fd7e14}
\makeatother
\makeatletter
\@ifpackageloaded{bookmark}{}{\usepackage{bookmark}}
\makeatother
\makeatletter
\@ifpackageloaded{caption}{}{\usepackage{caption}}
\AtBeginDocument{%
\ifdefined\contentsname
  \renewcommand*\contentsname{Table of contents}
\else
  \newcommand\contentsname{Table of contents}
\fi
\ifdefined\listfigurename
  \renewcommand*\listfigurename{List of Figures}
\else
  \newcommand\listfigurename{List of Figures}
\fi
\ifdefined\listtablename
  \renewcommand*\listtablename{List of Tables}
\else
  \newcommand\listtablename{List of Tables}
\fi
\ifdefined\figurename
  \renewcommand*\figurename{Figure}
\else
  \newcommand\figurename{Figure}
\fi
\ifdefined\tablename
  \renewcommand*\tablename{Table}
\else
  \newcommand\tablename{Table}
\fi
}
\@ifpackageloaded{float}{}{\usepackage{float}}
\floatstyle{ruled}
\@ifundefined{c@chapter}{\newfloat{codelisting}{h}{lop}}{\newfloat{codelisting}{h}{lop}[chapter]}
\floatname{codelisting}{Listing}
\newcommand*\listoflistings{\listof{codelisting}{List of Listings}}
\makeatother
\makeatletter
\makeatother
\makeatletter
\@ifpackageloaded{caption}{}{\usepackage{caption}}
\@ifpackageloaded{subcaption}{}{\usepackage{subcaption}}
\makeatother
\usepackage{bookmark}
\IfFileExists{xurl.sty}{\usepackage{xurl}}{} % add URL line breaks if available
\urlstyle{same}
\hypersetup{
  pdftitle={rfaR: Reservoir Frequency Analysis in R},
  pdfauthor={Dan McGraw, P.E.; Andrew Verdin, PhD, P.H.; C. Haden Smith, P.E.},
  colorlinks=true,
  linkcolor={blue},
  filecolor={Maroon},
  citecolor={Blue},
  urlcolor={Blue},
  pdfcreator={LaTeX via pandoc}}


\title{rfaR: Reservoir Frequency Analysis in R}
\usepackage{etoolbox}
\makeatletter
\providecommand{\subtitle}[1]{% add subtitle to \maketitle
  \apptocmd{\@title}{\par {\large #1 \par}}{}{}
}
\makeatother
\subtitle{User and Technical Manual}
\author{Dan McGraw, P.E. \and Andrew Verdin, PhD, P.H. \and C. Haden
Smith, P.E.}
\date{2026-11-13}
\begin{document}
\maketitle

\renewcommand*\contentsname{Table of contents}
{
\setcounter{tocdepth}{2}
\tableofcontents
}

\bookmarksetup{startatroot}

\chapter*{Preface}\label{preface}
\addcontentsline{toc}{chapter}{Preface}

\markboth{Preface}{Preface}

\texttt{rfaR} implements the RMC Reservoir Frequency Analysis (RMC-RFA)
methodology in R, producing reservoir stage-frequency curves with
uncertainty bounds by combining deterministic Modified Puls flood
routing with a nested Monte Carlo framework.

This manual complements the
\href{https://usace-rmc.github.io/rfaR/}{function reference site}
generated by pkgdown. Where the pkgdown site documents individual
functions, this manual explains the methodology, workflow, and how the
pieces fit together.

\section*{How to cite}\label{how-to-cite}
\addcontentsline{toc}{section}{How to cite}

\markright{How to cite}

\emph{Cover image: Stoner Creek Dam and USGS streamgage 03251900 near
Paris, Kentucky. Source: National Oceanic and Atmospheric Administration
(2026).}

\bookmarksetup{startatroot}

\chapter{Introduction}\label{introduction}

\texttt{rfaR} is an R implementation of the Risk Management Center
Reservoir Frequency Analysis (RMC-RFA) software. This package produces
reservoir hydrologic hazard curves with uncertainty bounds by combining
deterministic flood routing with a nested Monte Carlo framework.
Hydrologic hazard analyses at dam and levee sites support risk
assessments in accordance with the USACE Dam \& Levee Safety Program
(DLS).

\section{Purpose}\label{purpose}

The intended use of \texttt{rfaR} is to generate hydrologic hazard
curves/loading curves\footnote{In risk analysis, the \emph{hazard} is
  the risk source where potential consequences relate to harm (Aven et
  al., n.d.). In the context of dam and levee risk analyses, the source
  of risk is often defined using a reservoir stage-frequency
  relationship. The terms \emph{hydrologic loading},
  \emph{stage-frequency}, and \emph{hydrologic hazard} are used
  interchangeably to relate reservoir/river stages to annual exceedance
  probabilities (AEPs). These relationships are commonly presented as
  curves and are used as probabilistic loading conditions in dam and
  levee risk analyses.}, supporting analysis for dam and levee risk
analyses. The development of rfaR was primarily to conduct the nested
monte carlo frame work using GEV inflow volume-frequency parameters
developed from partial duration series in arid locations (RMC-RFA only
supports LPIII parameters). Additional development allows for analysis
at levees (no reservoir routing), transforming sampled unregulated
inflow values to regulated inflow values, and a stage-based
probabilistic bounding of hydrologic hazard curves.

\section{rfaR Analysis Summary}\label{rfar-analysis-summary}

\texttt{rfaR} combines deterministic flood routing (Modified Puls) with
a nested, two-loop Monte Carlo framework to quantify:

\begin{itemize}
\item
  \textbf{Natural variability} --- simulated in the inner loop as a
  \emph{realization}. Each realization represents 10,000 indiviual flood
  events with randomly sampled starting stages, hydrograph shapes, and
  inflow volumes (from an associated inflow-frequency distribution).
\item
  \textbf{Knowledge uncertainty (epistemic uncertainty)} --- simulated
  in the outer loop, by sampling different inflow-frequency parameter
  sets (from LP3 or GEV flood frequency distributions, informed by
  RMC-BestFit's MCMC output).
\end{itemize}

\section{Scope and Limitations}\label{scope-and-limitations}

\texttt{rfaR} reimplements the core RMC-RFA methodology in open-source
R. \texttt{rfaR} is intended to serve as an interim capability between
the current release of RMC-RFA and the planned second version of
RMC-RFA. As of November 2026, users cannot import an existing RMC-RFA
model or input file directly into \texttt{rfaR}; users must prepare the
required inputs separately.

\section{How This Manual Is
Organized}\label{how-this-manual-is-organized}

\begin{itemize}
\tightlist
\item
  \textbf{Installation} --- System requirements and instructions for
  installing and loading \texttt{rfaR}.
\item
  \textbf{Set Up an rfaR Analysis} --- Prepare a new \texttt{rfaR}
  analysis, including project organization and input-data preparation.
\item
  \textbf{Methodology} --- Statistical and hydrologic methods underlying
  the two-loop simulation framework.
\item
  \textbf{Workflow} --- Detailed guidance for using
  \texttt{rfa\_simulate()}, defining input options and simulation
  settings, and processing outputs.
\item
  \textbf{Validation} --- Compare \texttt{rfaR} results with RMC-RFA and
  HEC-HMS benchmark results.
\item
  \textbf{Function Reference} --- Reference documentation for
  \texttt{rfaR} functions and their arguments.
\end{itemize}

For function-level documentation, also see the
\href{https://usace-rmc.github.io/rfaR/}{pkgdown reference site}.

\bookmarksetup{startatroot}

\chapter{Installation}\label{installation}

\texttt{rfaR} is not currently available on the Comprehensive R Archive
Network (CRAN). The package source code is available from the
\href{https://github.com/USACE-RMC/rfaR}{\texttt{USACE-RMC/rfaR} GitHub
repository}.

\begin{tcolorbox}[enhanced jigsaw, arc=.35mm, bottomrule=.15mm, bottomtitle=1mm, breakable, colback=white, colbacktitle=quarto-callout-warning-color!10!white, colframe=quarto-callout-warning-color-frame, coltitle=black, left=2mm, leftrule=.75mm, opacityback=0, opacitybacktitle=0.6, rightrule=.15mm, title=\textcolor{quarto-callout-warning-color}{\faExclamationTriangle}\hspace{0.5em}{Warning}, titlerule=0mm, toprule=.15mm, toptitle=1mm]

\texttt{rfaR} is under active development. Its application programming
interface (API) may change without notice.

\end{tcolorbox}

\section{System Requirements}\label{system-requirements}

\begin{itemize}
\tightlist
\item
  R version 4.1 or later
\item
  access to the \texttt{USACE-RMC/rfaR} GitHub repository
\item
  Recommended: RStudio
  \href{https://posit.co/products/open-source/rstudio}{RStudio} or
  \href{https://positron.posit.co/}{Positron}
\end{itemize}

Installing \texttt{rfaR} with \texttt{remotes} or \texttt{devtools}
installs required R package dependencies as needed.

\section{Install rfaR}\label{install-rfar}

\subsection{Install with remotes}\label{install-with-remotes}

Install the \texttt{remotes} package, if needed, and then install
\texttt{rfaR} from GitHub:

\begin{Shaded}
\begin{Highlighting}[]
\FunctionTok{install.packages}\NormalTok{(}\StringTok{"remotes"}\NormalTok{)}

\NormalTok{remotes}\SpecialCharTok{::}\FunctionTok{install\_github}\NormalTok{(}
  \StringTok{"USACE{-}RMC/rfaR"}\NormalTok{,}
  \AttributeTok{build\_vignettes =} \ConstantTok{FALSE}
\NormalTok{)}
\end{Highlighting}
\end{Shaded}

\subsection{Developers and
Contributors}\label{developers-and-contributors}

Users who plan to modify the package source should clone the repository
and work from a local copy of the project. After cloning the repository
with Git or GitHub Desktop, open an R session or RStudio Project from
the repository root directory.

Install devtools if it is not already available:

\begin{Shaded}
\begin{Highlighting}[]
\FunctionTok{install.packages}\NormalTok{(}\StringTok{"devtools"}\NormalTok{)}

\NormalTok{devtools}\SpecialCharTok{::}\FunctionTok{install\_github}\NormalTok{(}\StringTok{"USACE{-}RMC/rfaR"}\NormalTok{)}
\end{Highlighting}
\end{Shaded}

During active development, load the package source into the current R
session without installing it:

\begin{Shaded}
\begin{Highlighting}[]
\NormalTok{devtools}\SpecialCharTok{::}\FunctionTok{load\_all}\NormalTok{()}
\end{Highlighting}
\end{Shaded}

To install the local working copy of \texttt{rfaR} into your R library,
run:

\begin{Shaded}
\begin{Highlighting}[]
\NormalTok{devtools}\SpecialCharTok{::}\FunctionTok{install}\NormalTok{()}
\end{Highlighting}
\end{Shaded}

\begin{tcolorbox}[enhanced jigsaw, arc=.35mm, bottomrule=.15mm, bottomtitle=1mm, breakable, colback=white, colbacktitle=quarto-callout-note-color!10!white, colframe=quarto-callout-note-color-frame, coltitle=black, left=2mm, leftrule=.75mm, opacityback=0, opacitybacktitle=0.6, rightrule=.15mm, title=\textcolor{quarto-callout-note-color}{\faInfo}\hspace{0.5em}{Note}, titlerule=0mm, toprule=.15mm, toptitle=1mm]

Cloning the repository does not make \texttt{library(rfaR)} available.
To use the package, install it with \texttt{devtools::install()} or
\texttt{remotes::install\_github()}, or load the source for development
with \texttt{devtools::load\_all()}.

\end{tcolorbox}

\section{Dependencies}\label{dependencies}

Required package dependencies are declared in the \texttt{rfaR}
\texttt{DESCRIPTION} file and are installed automatically during
installation.

USACE users may encounter package download, installation, or
compatibility issues because of network controls, repository-access
restrictions, or locally managed R environments.

To report an installation or version-compatibility issue, open an issue
in the \href{https://github.com/USACE-RMC/rfaR/issues}{\texttt{rfaR}
GitHub repository}. Include the R version, operating system,
installation method, and complete error message when possible. Do not
include Controlled Unclassified Information (CUI), credentials, access
tokens, or other nonpublic system information in a GitHub issue.

\section{Verify Installation}\label{verify-installation}

Load \texttt{rfaR} and confirm that the installed version and included
example data are available:

\begin{Shaded}
\begin{Highlighting}[]
\FunctionTok{library}\NormalTok{(rfaR)}

\FunctionTok{packageVersion}\NormalTok{(}\StringTok{"rfaR"}\NormalTok{)}

\FunctionTok{head}\NormalTok{(jmd\_wy1980\_stage)}
\end{Highlighting}
\end{Shaded}

If these commands run without error, rfaR is installed and ready to use.

\bookmarksetup{startatroot}

\chapter{Set Up an rfaR Analysis}\label{set-up-an-rfar-analysis}

This chapter will step through best practices to set up, create, and
execute an \texttt{rfaR} analysis.

\section{Create a Project Directory}\label{create-a-project-directory}

Create a dedicated directory for each \texttt{rfaR} analysis. Store
project inputs, scripts, and outputs in a location where you have write
access, such as your user Documents directory. Do not create analysis
files in protected system directories, such as
\texttt{C:\textbackslash{}Program\ Files}.

The following example creates an \texttt{rfaR-analysis} directory in the
current user's home directory and creates subdirectories for input data,
scripts, and outputs. The commands will create a folder structure to
support a single rfaR analysis (R \& unix users use the term
\emph{directory} as opposed to \emph{folder}). Csers can create the same
folder structure using Windows Explorer.

\begin{Shaded}
\begin{Highlighting}[]
\CommentTok{\# Set a base directory}
\CommentTok{\# base\_dir \textless{}{-} file.path("C:/path/to/your/project/")}

\CommentTok{\# Define the project directory (replace path.expand("\textasciitilde{}") with base\_dir)}
\NormalTok{project\_dir }\OtherTok{\textless{}{-}} \FunctionTok{file.path}\NormalTok{(}\FunctionTok{path.expand}\NormalTok{(}\StringTok{"\textasciitilde{}"}\NormalTok{), }\StringTok{"rfaR{-}analysis"}\NormalTok{)}

\CommentTok{\# Define project subdirectories}
\NormalTok{project\_dirs }\OtherTok{\textless{}{-}} \FunctionTok{c}\NormalTok{(}
\NormalTok{  project\_dir,}
  \FunctionTok{file.path}\NormalTok{(project\_dir, }\StringTok{"data"}\NormalTok{),}
  \FunctionTok{file.path}\NormalTok{(project\_dir, }\StringTok{"data"}\NormalTok{, }\StringTok{"raw"}\NormalTok{),}
  \FunctionTok{file.path}\NormalTok{(project\_dir, }\StringTok{"data"}\NormalTok{, }\StringTok{"processed"}\NormalTok{),}
  \FunctionTok{file.path}\NormalTok{(project\_dir, }\StringTok{"output"}\NormalTok{),}
  \FunctionTok{file.path}\NormalTok{(project\_dir, }\StringTok{"scripts"}\NormalTok{)}
\NormalTok{)}

\CommentTok{\# Create directories if they do not already exist}
\FunctionTok{invisible}\NormalTok{(}\FunctionTok{lapply}\NormalTok{(}
\NormalTok{  project\_dirs,}
\NormalTok{  dir.create,}
  \AttributeTok{recursive =} \ConstantTok{TRUE}\NormalTok{,}
  \AttributeTok{showWarnings =} \ConstantTok{FALSE}
\NormalTok{))}
\end{Highlighting}
\end{Shaded}

The resulting directory structure is now:

\begin{verbatim}
rfaR-analysis/
├── data/
│   ├── raw/
│   └── processed/
├── output/
└── scripts/
\end{verbatim}

\begin{itemize}
\tightlist
\item
  Store original source data in \texttt{data/raw/}.
\item
  Store cleaned, reformatted, or derived datasets in
  \texttt{data/processed/}.
\item
  Store analysis scripts in \texttt{scripts/}.
\item
  Store simulation results, figures, tables, and exported files in
  \texttt{output/}.
\end{itemize}

\section{\texorpdfstring{Import Data from Existing RMC-RFA Projects
using
\texttt{rfaScrape}}{Import Data from Existing RMC-RFA Projects using rfaScrape}}\label{import-data-from-existing-rmc-rfa-projects-using-rfascrape}

Direct integration between RMC-RFA SQLite project files and
\texttt{rfaR} is under development. Until that capability is available,
the recommended approach for extracting organized input data and
analysis results from an existing RMC-RFA project is to use the
\href{https://github.com/danmcgraw94/rfaScrape}{\texttt{rfaScrape}}
package.

\texttt{rfaScrape} extracts data and results from RMC-RFA project files
with the \texttt{.rfa.sqlite} extension and exports them to organized
CSV directories. The package is intended to improve the organization,
transfer, and reporting of RMC-RFA project data for use in downstream
analyses.

Future \texttt{rfaR} development is expected to incorporate
functionality that supports more direct use of data from existing
RMC-RFA projects.

\subsection{R Project}\label{r-project}

Users working in RStudio may optionally create an RStudio Project in the
\texttt{rfaR-analysis} directory. An RStudio Project helps maintain
consistent working directory but is not required to run \texttt{rfaR}.

\section{Load Input Data}\label{load-input-data}

\texttt{rfaR} analyses currently use input data stored in
comma-separated values (CSV) files. Export or save the required input
datasets as CSV files and store them in the project \texttt{data/raw/}
directory. This manual uses the base R function \texttt{read.csv()} for
all examples. Users may instead use \texttt{readr::read\_csv()} if they
prefer the \texttt{readr} import workflow; however, the resulting data
must meet the formatting requirements described in the following
sections.

The example below imports a CSV file from the project input-data
directory:

\begin{Shaded}
\begin{Highlighting}[]
\CommentTok{\# Define the path to an input CSV file}
\NormalTok{input\_file }\OtherTok{\textless{}{-}} \FunctionTok{file.path}\NormalTok{(}
\NormalTok{  project\_dir,}
  \StringTok{"data"}\NormalTok{,}
  \StringTok{"raw"}\NormalTok{,}
  \StringTok{"example\_input.csv"}
\NormalTok{)}

\CommentTok{\# Import the CSV file}
\NormalTok{input\_data }\OtherTok{\textless{}{-}} \FunctionTok{read.csv}\NormalTok{(input\_file)}

\CommentTok{\# Review the imported data}
\FunctionTok{head}\NormalTok{(input\_data)}
\FunctionTok{str}\NormalTok{(input\_data)}
\end{Highlighting}
\end{Shaded}

Use \texttt{head()} to review the first rows of the imported dataset.
Use \texttt{str()} to confirm the number of observations, variable
names, and R data types.

Users who are new to importing, cleaning, or restructuring data in R may
find the following resources helpful:

\begin{itemize}
\tightlist
\item
  \href{https://stat.ethz.ch/R-manual/R-devel/library/utils/html/read.table.html}{\texttt{read.csv()}}(\emph{\texttt{read.table()}})
\item
  \href{https://r4ds.hadley.nz/data-import.html}{Data Import}
\item
  \href{https://r4ds.hadley.nz/data-transform.html}{Data Transform}
\item
  \href{https://readr.tidyverse.org/reference/read_delim.html}{\texttt{read\_csv()}}
\end{itemize}

\section{Input Data}\label{input-data}

The input data and required formatting is included below. This section
will be updated to reflect on-going package development

\subsection{Inflow Hydrographs}\label{inflow-hydrographs}

Inflow hydrographs define the temporal distribution, or shape, of flood
inflow events used during simulation. Each hydrograph represents a
historical, synthetic, design, or other representative flood event.
During a simulation, \texttt{rfaR} samples a hydrograph shape and scales
it to the sampled inflow magnitude.

Use \texttt{hydrograph\_setup()} to prepare one or more input
hydrographs for use in \texttt{rfa\_simulate()}. Input hydrographs may
be copied directly from RMC-RFA or imported from similarly formatted CSV
files.

\subsubsection{Required Format}\label{required-format}

Each input hydrograph must be a data frame with the following columns in
this order:

{\def\LTcaptype{none} % do not increment counter
\begin{longtable}[]{@{}
  >{\raggedright\arraybackslash}p{(\linewidth - 4\tabcolsep) * \real{0.3333}}
  >{\raggedright\arraybackslash}p{(\linewidth - 4\tabcolsep) * \real{0.3333}}
  >{\raggedright\arraybackslash}p{(\linewidth - 4\tabcolsep) * \real{0.3333}}@{}}
\toprule\noalign{}
\begin{minipage}[b]{\linewidth}\raggedright
Column
\end{minipage} & \begin{minipage}[b]{\linewidth}\raggedright
Description
\end{minipage} & \begin{minipage}[b]{\linewidth}\raggedright
Required format
\end{minipage} \\
\midrule\noalign{}
\endhead
\bottomrule\noalign{}
\endlastfoot
\texttt{Ordinate} & Sequential row or ordinate number. & Integer or
numeric. \\
\texttt{Date} & Date associated with each inflow observation. &
Character string in month/day/year format, such as
\texttt{"2/2/2099"}. \\
\texttt{Time} & Time associated with each inflow observation. &
Character string in hour:minute format, such as \texttt{"0:00"} or
\texttt{"13:30"}. \\
\texttt{Flow} & Inflow discharge for each time step. & Numeric flow in
cubic feet per second (cfs). Character values containing comma
separators are also supported. \\
\end{longtable}
}

The hydrograph records must be in chronological order and use a regular
time step. \texttt{hydrograph\_setup()} determines the time step from
the first two observations in each hydrograph.

The following example shows the format of an example hydrograph
(\texttt{jmd\_hydro\_pmf}):

\begin{Shaded}
\begin{Highlighting}[]
\NormalTok{  Ordinate     Date Time Flow}
\NormalTok{1        1 2/2/2099 0:00 1131}
\NormalTok{2        2 2/2/2099 1:00 1890}
\NormalTok{3        3 2/2/2099 2:00 2277}
\NormalTok{4        4 2/2/2099 3:00 1936}
\NormalTok{5        5 2/2/2099 4:00 1934}
\NormalTok{6        6 2/2/2099 5:00 1933}
\end{Highlighting}
\end{Shaded}

Although the \texttt{Ordinate} column is retained in the input format
for consistency with RMC-RFA exports, \texttt{rfaR} uses the
\texttt{Date}, \texttt{Time}, and \texttt{Flow} columns to construct the
simulation hydrograph.

\subsubsection{Import Hydrograph Data}\label{import-hydrograph-data}

Store hydrograph CSV files in the project \texttt{data/raw/} directory.
The following example imports three hydrographs:

\begin{Shaded}
\begin{Highlighting}[]
\NormalTok{hydro\_pmf }\OtherTok{\textless{}{-}} \FunctionTok{read.csv}\NormalTok{(}
  \FunctionTok{file.path}\NormalTok{(project\_dir, }\StringTok{"data"}\NormalTok{, }\StringTok{"raw"}\NormalTok{, }\StringTok{"hydro\_pmf.csv"}\NormalTok{),}
  \AttributeTok{stringsAsFactors =} \ConstantTok{FALSE}
\NormalTok{)}

\NormalTok{hydro\_jun1965 }\OtherTok{\textless{}{-}} \FunctionTok{read.csv}\NormalTok{(}
  \FunctionTok{file.path}\NormalTok{(project\_dir, }\StringTok{"data"}\NormalTok{, }\StringTok{"raw"}\NormalTok{, }\StringTok{"hydro\_jun1965.csv"}\NormalTok{),}
  \AttributeTok{stringsAsFactors =} \ConstantTok{FALSE}
\NormalTok{)}

\NormalTok{hydro\_may1955 }\OtherTok{\textless{}{-}} \FunctionTok{read.csv}\NormalTok{(}
  \FunctionTok{file.path}\NormalTok{(project\_dir, }\StringTok{"data"}\NormalTok{, }\StringTok{"raw"}\NormalTok{, }\StringTok{"hydro\_may1955.csv"}\NormalTok{),}
  \AttributeTok{stringsAsFactors =} \ConstantTok{FALSE}
\NormalTok{)}
\end{Highlighting}
\end{Shaded}

Review each imported data frame before continuing:

\begin{Shaded}
\begin{Highlighting}[]
\FunctionTok{head}\NormalTok{(hydro\_pmf)}
\FunctionTok{str}\NormalTok{(hydro\_pmf)}
\end{Highlighting}
\end{Shaded}

\subsubsection{Prepare Hydrographs for
Simulation}\label{prepare-hydrographs-for-simulation}

Use \texttt{hydrograph\_setup()} to format the imported hydrographs for
simulation.

\begin{Shaded}
\begin{Highlighting}[]
\NormalTok{hydrographs }\OtherTok{\textless{}{-}} \FunctionTok{hydrograph\_setup}\NormalTok{(}
\NormalTok{  hydro\_pmf,}
\NormalTok{  hydro\_jun1965,}
\NormalTok{  hydro\_may1955,}
  \AttributeTok{critical\_duration =} \DecValTok{2}\NormalTok{,}
  \AttributeTok{routing\_days =} \DecValTok{10}
\NormalTok{)}
\end{Highlighting}
\end{Shaded}

The \texttt{critical\_duration} argument specifies the inflow duration,
in days, used to calculate the critical-duration inflow statistic for
each hydrograph. The critical duration should align with the duration
used for the inflow-volume frequency analysis. It must also correspond
to a whole number of hydrograph time steps.

The \texttt{routing\_days} argument specifies the desired total
duration, in days, of each hydrograph used in routing. If a hydrograph
is shorter than the specified routing duration,
\texttt{hydrograph\_setup()} extends the hydrograph by appending
zero-flow time steps.

\subsubsection{Hydrograph Sampling
Weights}\label{hydrograph-sampling-weights}

By default, \texttt{rfaR} assigns equal sampling probability to each
supplied hydrograph. Use the optional \texttt{weights} argument to
assign different relative sampling weights.

\begin{Shaded}
\begin{Highlighting}[]
\NormalTok{hydrographs }\OtherTok{\textless{}{-}} \FunctionTok{hydrograph\_setup}\NormalTok{(}
\NormalTok{  hydro\_pmf,}
\NormalTok{  hydro\_jun1965,}
\NormalTok{  hydro\_may1955,}
  \AttributeTok{critical\_duration =} \DecValTok{2}\NormalTok{,}
  \AttributeTok{routing\_days =} \DecValTok{10}\NormalTok{,}
  \AttributeTok{weights =} \FunctionTok{c}\NormalTok{(}\DecValTok{3}\NormalTok{, }\DecValTok{1}\NormalTok{, }\DecValTok{1}\NormalTok{)}
\NormalTok{)}
\end{Highlighting}
\end{Shaded}

In this example, \texttt{hydro\_pmf} has three times the relative
sampling weight of each historical hydrograph. The supplied weights are
normalized internally to produce sampling probabilities.

View the normalized probabilities with:

\begin{Shaded}
\begin{Highlighting}[]
\FunctionTok{attr}\NormalTok{(hydrographs, }\StringTok{"probs"}\NormalTok{)}
\end{Highlighting}
\end{Shaded}

\subsubsection{Output}\label{output}

\texttt{hydrograph\_setup()} returns a list of formatted hydrograph data
frames. Each formatted hydrograph includes the following variables:

{\def\LTcaptype{none} % do not increment counter
\begin{longtable}[]{@{}
  >{\raggedright\arraybackslash}p{(\linewidth - 2\tabcolsep) * \real{0.5000}}
  >{\raggedright\arraybackslash}p{(\linewidth - 2\tabcolsep) * \real{0.5000}}@{}}
\toprule\noalign{}
\begin{minipage}[b]{\linewidth}\raggedright
Variable
\end{minipage} & \begin{minipage}[b]{\linewidth}\raggedright
Description
\end{minipage} \\
\midrule\noalign{}
\endhead
\bottomrule\noalign{}
\endlastfoot
\texttt{datetime} & Date-time value converted to \texttt{POSIXct}. \\
\texttt{hour} & Elapsed hours from the start of the hydrograph. \\
\texttt{inflow} & Inflow discharge in cfs. \\
\texttt{hydrograph\_num} & Numeric identifier used during hydrograph
sampling. \\
\end{longtable}
}

Each hydrograph also stores the following attributes:

{\def\LTcaptype{none} % do not increment counter
\begin{longtable}[]{@{}
  >{\raggedright\arraybackslash}p{(\linewidth - 2\tabcolsep) * \real{0.5000}}
  >{\raggedright\arraybackslash}p{(\linewidth - 2\tabcolsep) * \real{0.5000}}@{}}
\toprule\noalign{}
\begin{minipage}[b]{\linewidth}\raggedright
Attribute
\end{minipage} & \begin{minipage}[b]{\linewidth}\raggedright
Description
\end{minipage} \\
\midrule\noalign{}
\endhead
\bottomrule\noalign{}
\endlastfoot
\texttt{obs\_vol} & Critical-duration inflow statistic calculated from
the observed hydrograph. \\
\texttt{dt} & Hydrograph time step in hours. \\
\texttt{name} & Name of the input hydrograph object. \\
\end{longtable}
}

The returned hydrograph list stores normalized sampling probabilities in
its \texttt{probs} attribute.

\begin{Shaded}
\begin{Highlighting}[]
\CommentTok{\# Review the first formatted hydrograph}
\FunctionTok{head}\NormalTok{(hydrographs[[}\DecValTok{1}\NormalTok{]])}

\CommentTok{\# Review hydrograph attributes}
\FunctionTok{attributes}\NormalTok{(hydrographs[[}\DecValTok{1}\NormalTok{]])}

\CommentTok{\# Review hydrograph sampling probabilities}
\FunctionTok{attr}\NormalTok{(hydrographs, }\StringTok{"probs"}\NormalTok{)}
\end{Highlighting}
\end{Shaded}

\begin{quote}
\textbf{Note:} \texttt{hydrograph\_setup()} removes trailing rows with
missing flow values, which can occur when CSV files are exported from
spreadsheet software. Review the input data before analysis and correct
missing or malformed values within the hydrograph record.
\end{quote}

\subsection{Bestfit Parameter Sets}\label{bestfit-parameter-sets}

\subsection{Flood Seasonality}\label{flood-seasonality}

\subsection{Observed Stage Timeseries}\label{observed-stage-timeseries}

\bookmarksetup{startatroot}

\chapter{Methodology}\label{methodology}

\section{Overview}\label{overview}

\texttt{rfaR} implements a two-loop, nested Monte Carlo framework to
propagate both natural variability and knowledge uncertainty through a
deterministic flood routing model.

\section{Outer Loop: Parameter
Uncertainty}\label{outer-loop-parameter-uncertainty}

\section{Inner Loop: Realizations}\label{inner-loop-realizations}

\section{Flood Routing}\label{flood-routing}

\bookmarksetup{startatroot}

\chapter{Workflow}\label{workflow}

\section{Overview}\label{overview-1}

\section{Step 1: Prepare Inputs}\label{step-1-prepare-inputs}

\section{Step 2: Hydrograph Setup}\label{step-2-hydrograph-setup}

\section{Step 3: Sampling}\label{step-3-sampling}

\section{Step 4: Routing}\label{step-4-routing}

\section{Step 5: Running the Full
Simulation}\label{step-5-running-the-full-simulation}

\section{Step 6: Interpreting
Results}\label{step-6-interpreting-results}

\bookmarksetup{startatroot}

\chapter{Validation}\label{validation}

\section{Purpose}\label{purpose-1}

\section{Modified Puls Routing}\label{modified-puls-routing}

\section{Stratified Sampling}\label{stratified-sampling}

\section{See Also}\label{see-also}

Detailed validation vignettes are also available on the
\href{https://usace-rmc.github.io/rfaR/articles/index.html}{pkgdown
site}.

\bookmarksetup{startatroot}

\chapter{Function Reference}\label{function-reference}

\section{Execute Reservoir Frequency
Analysis}\label{execute-reservoir-frequency-analysis}

\subsection{\texorpdfstring{\texttt{rfa\_simulate()}}{rfa\_simulate()}}\label{rfa_simulate}

\begin{verbatim}
RFA Stage-Frequency Simulation

Description:

     Performs reservoir stage frequency analysis using Monte Carlo
     simulation to develop stage-frequency relationships. Combines
     stratified flow-frequency sampling, flood seasonality, hydrograph
     scaling, and Modified Puls reservoir routing to estimate annual
     exceedance probabilities of reservoir stage.

Usage:

     rfa_simulate(
       sim_type = "expected",
       bestfit_params,
       dist = "LP3",
       stage_ts,
       seasonality,
       hydrographs,
       resmodel,
       Nbins = 50,
       events_per_bin = 200,
       routing_dt = 1,
       Ncores = NULL,
       sim_name = NULL,
       results_dir = NULL
     )
     
Arguments:

sim_type: Character string specifying the simulation mode. One of
          ‘"median"’, ‘"expected"’, or ‘"full"’. Default is
          ‘"expected"’.

bestfit_params: Data frame or matrix of distribution parameters from
          RMC-BestFit MCMC output. Columns 1-3 are distribution
          parameters:

            * LP3: mean (log), standard deviation (log), skew (log)

            * GEV: location, scale, shape

          Column 4 is log-likelihood (used only in ‘"median"’ mode to
          identify the posterior mode). For ‘"expected"’ mode, must
          have at least ‘Nbins * events_per_bin’ rows. For ‘"full"’
          mode, each row is treated as an independent realization.

    dist: Character string specifying the frequency distribution.
          Either ‘"LP3"’ (Log-Pearson Type III, default) or ‘"GEV"’
          (Generalized Extreme Value).

stage_ts: Data frame of historical reservoir stage with columns ‘date’
          (character in M/D/YYYY format) and ‘stage’ (numeric, feet).

seasonality: Numeric vector of length 12 giving monthly flood
          occurrence probabilities (relative frequencies). Used to
          sample the month of each simulated event, which determines
          the antecedent pool elevation.

hydrographs: List of hydrograph data frames as returned by
          ‘hydrograph_setup’. Each element has columns ‘datetime’,
          ‘hour’, ‘inflow’, and ‘hydrograph_num’, with attributes
          ‘"obs_vol"’ (observed max n-day volume) and ‘"dt"’ (timestep
          in hours). The list must have a ‘"probs"’ attribute
          containing normalized sampling probabilities.

resmodel: Data frame with three columns: elevation (ft), storage
          (acre-ft), and discharge (cfs). Defines the reservoir model
          for Modified Puls routing.

   Nbins: Integer. Number of stratified sampling bins in EV1 space.
          Default is 50.

events_per_bin: Integer. Number of events sampled per bin. Default is
          200. Total simulations per realization = ‘Nbins *
          events_per_bin’.

routing_dt: Numeric. Routing timestep in hours. Passed to
          ‘scale_hydrograph’ to resample hydrographs before routing.
          Supported values are ‘0.25’ (15-min), ‘1’ (1-hour, default),
          ‘6’ (6-hour), and ‘24’ (24-hour).

  Ncores: Integer or ‘NULL’. Number of parallel workers for ‘"full"’
          mode. If ‘NULL’ (default), automatically selects based on
          available cores (capped at 16). Ignored for ‘"median"’ and
          ‘"expected"’ modes.

sim_name: Optional character string to label the simulation. Used in
          the output CSV filename. If ‘NULL’ (default), defaults to
          ‘"sim"’. Spaces are replaced with underscores in the
          filename.

results_dir: Optional path to the directory where results are saved. If
          ‘NULL’ (default), creates an ‘rfaR_results’ folder in the
          current working directory.

Details:

     Three simulation modes are available:

     ‘"median"’ Single realization using the most likely (posterior
          mode) parameter set from RMC-BestFit. Produces one
          stage-frequency curve with no uncertainty bounds. Useful for
          quick screening or debugging.

     ‘"expected"’ Single realization using expected value sampling,
          where each stratified flow sample is paired 1-to-1 with a
          different parameter set from the posterior distribution.
          Produces one stage-frequency curve representing the expected
          (mean) result integrated across parameter uncertainty.

     ‘"full"’ Nested Monte Carlo with full uncertainty quantification.
          The outer loop iterates over all parameter sets (one per row
          of ‘bestfit_params’), and each inner loop independently
          samples natural variability via stratified sampling,
          seasonality, starting pool, and hydrograph shape. Each
          realization produces an independent stage-frequency curve;
          curves are combined on a common stage grid to yield expected,
          median, and 5th/95th percentile confidence bounds. The outer
          loop is parallelized using the ‘future’ framework.

     For all modes, natural variability is captured through stratified
     sampling in with ‘Nbins’ bins and ‘events_per_bin’ events per bin,
     providing reliable coverage from common events down to
     approximately 1e-8 AEP.

     The stratified sampling approach divides the AEP range (default
     0.99 to 1e-8) into equal-width bins in EV1 (Gumbel reduced
     variate) space and samples events uniformly within each bin. This
     ensures adequate representation of rare flood events that would
     require orders of magnitude more samples under crude Monte Carlo
     sampling.

     Post-processing uses the law of total probability via
     ‘stage_frequency_curve’: for a grid of stage thresholds, the
     weighted exceedance fraction is computed within each bin and
     summed across bins.

     In ‘"full"’ mode, each parallel worker independently samples all
     random inputs (seasonality, starting pool, hydrograph shape, and
     stratified z-variates) in addition to using its own parameter set.
     This ensures full Monte Carlo independence across realizations.
     Each worker post-processes its results into a stage-frequency
     curve immediately, returning only the curve rather than the full
     peak stage matrix, keeping memory usage manageable for large
     numbers of realizations.

     Curves from individual realizations are combined by interpolating
     each onto a common stage grid (spanning the global min/max across
     all realizations) in log10(AEP) space, then computing summary
     statistics across realizations at each stage threshold.

     Results are automatically exported as a CSV file named
     ‘{sim_name}_{sim_type}_{MM_DD_YY_HHMM}.csv’ in the ‘results_dir’
     directory. For example, a median simulation named "John McGraw
     Dam" run on June 4, 2025 at 2:30 PM would produce
     ‘John_McGraw_Dam_median_06_04_25_1430.csv’.

Value:

     A list whose contents depend on ‘sim_type’:

     For ‘"median"’ and ‘"expected"’:

     stage_frequency Data frame with columns ‘stage’ and ‘AEP’ from
          ‘stage_frequency_curve’.

     peakStage Matrix of simulated peak stages (‘events_per_bin’ rows
          by ‘Nbins’ columns).

     peakFlow Matrix of simulated peak discharges (same dimensions).

     weights Stratified sampling bin weights.

     For ‘"full"’:

     stage_frequency Data frame with columns ‘stage’, ‘expected’,
          ‘median’, ‘lower_05’, and ‘upper_95’ on a common stage grid.

     all_curves List of per-realization stage-frequency data frames
          (each with ‘stage’ and ‘AEP’ columns).

     aep_matrix Matrix of interpolated AEP values (500 rows by
          ‘Nrealizations’ columns) on the common stage grid.

     common_stage Numeric vector of the common stage grid (length 500).

     Nrealizations Number of parameter set realizations processed.

     Nsims_per_realiz Number of simulations per realization.

See Also:

     ‘flow_frequency_sampler’, ‘flow_frequency_sampler_expected’,
     ‘stratified_sampler’, ‘scale_hydrograph’, ‘mod_puls_routing’,
     ‘stage_frequency_curve’, ‘hydrograph_setup’

Examples:

     ## Not run:
     
     # Not run: Expected and Median modes take ~15s; Full uncert
     # runs ~10,000 parameter realizations and takes 3-4 hours on a
     # typical machine.
     
     # --- Setup ---
     hydros <- hydrograph_setup(jmd_hydro_apr1999, jmd_hydro_jun1965,
                                jmd_hydro_may1955, jmd_hydro_pmf,
                                critical_duration = 2, routing_days = 10)
     
     # --- Expected only (default) ---
     results_exp <- rfa_simulate(
       sim_type       = "expected",
       bestfit_params = jmd_bf_parameter_sets,
       stage_ts       = jmd_wy1980_stage,
       seasonality    = jmd_seasonality$relative_frequency,
       hydrographs    = hydros,
       resmodel       = jmd_resmodel,
       sim_name       = "jmd"
     )
     
     # --- Median only ---
     results_med <- rfa_simulate(
       sim_type       = "median",
       bestfit_params = jmd_bf_parameter_sets,
       stage_ts       = jmd_wy1980_stage,
       seasonality    = jmd_seasonality$relative_frequency,
       hydrographs    = hydros,
       resmodel       = jmd_resmodel,
       sim_name       = "jmd"
     )
     
     # --- Full uncertainty (parallelized) ---
     results_full <- rfa_simulate(
       sim_type       = "full",
       bestfit_params = jmd_bf_parameter_sets,
       stage_ts       = jmd_wy1980_stage,
       seasonality    = jmd_seasonality$relative_frequency,
       hydrographs    = hydros,
       resmodel       = jmd_resmodel,
       Ncores         = 4,
       sim_name       = "jmd"
     )
     ## End(Not run)
     
\end{verbatim}

\section{Primary RFA Modules}\label{primary-rfa-modules}

\subsection{\texorpdfstring{\texttt{stratified\_sampler()}}{stratified\_sampler()}}\label{stratified_sampler}

\begin{verbatim}
Stratified Sampler

Description:

     Creates stratified bins in standard normal space for use in
     stratified sampling. Stratification improves sampling efficiency
     by ensuring adequate coverage of rare events.

Usage:

     stratified_sampler(
       minAEP = 1e-08,
       maxAEP = 0.99,
       dist = "EV1",
       Nbins = NULL,
       Mevents = NULL,
       verbose = FALSE
     )
     
Arguments:

  minAEP: Minimum annual exceedance probability. Default is ‘1E-8’.

  maxAEP: Maximum annual exceedance probability. Default is ‘0.99’.

    dist: Character. Probability space for stratification.
          Case-insensitive. Default is ‘"EV1"’. Invalid values trigger
          a warning and default to ‘"EV1"’.

          ‘"EV1"’ Extreme Value Type I (Gumbel) space. Recommended for
              flood frequency analysis. Allocates more bins to rare
              events through the transformation ‘-log(-log(1-AEP))’,
              improving tail estimation efficiency.

          ‘"Normal"’ Standard normal (z-score) space. Uniform bins in
              z-space. Use when the underlying phenomenon is normally
              distributed.

          ‘"Uniform"’ Uniform probability space. Equal probability
              width per bin. Generally inefficient for rare event
              estimation.

   Nbins: Number of stratified bins. Default is ‘20’.

 Mevents: Number of events per bin. Default is ‘500’.

 verbose: Logical. If ‘TRUE’, prints a completion message summarizing
          the stratification. Default is ‘FALSE’.

Details:

     The function divides the probability space into bins using the
     specified transformation. EV1 transformation is recommended for
     heavy-tailed distributions common in flood frequency analysis, as
     it naturally allocates more sampling effort to rare events
     critical for dam safety assessments.

Value:

     A list containing:

     normOrd Vector of standard normal ordinates spanning all bins

     Weights Vector of probability weights for each bin

     Zlower Vector of lower bounds for each bin (standard normal)

     Zupper Vector of upper bounds for each bin (standard normal)

     Nbins Number of bins

     Mevents Number of events per bin

Examples:

     # Default stratification
     strat <- stratified_sampler()
     
     # Custom bins and events
     strat <- stratified_sampler(minAEP = 1E-6,
                                 maxAEP = 0.5,
                                 dist = "EV1",
                                 Nbins = 10,
                                 Mevents = 100)
     
     # With verbose message
     strat <- stratified_sampler(verbose = TRUE)
     
\end{verbatim}

\subsection{\texorpdfstring{\texttt{flow\_frequency\_sampler()}}{flow\_frequency\_sampler()}}\label{flow_frequency_sampler}

\begin{verbatim}
Flow Frequency Sampler

Description:

     Generates a stratified matrix of flow values from a single set of
     frequency distribution parameters using stratified Monte Carlo
     sampling. Used internally by ‘rfa_simulate()’ for median-only and
     full uncertainty modes.

Usage:

     flow_frequency_sampler(
       bestfit_params,
       freq_dist = "LP3",
       strat_dist = "ev1",
       Nbin = NULL,
       Mevent = NULL
     )
     
Arguments:

bestfit_params: Numeric vector of length 3 containing distribution
          parameters. For LP3: ‘c(mean_log, sd_log, skew_log)’. For
          GEV: ‘c(location, scale, shape)’.

freq_dist: Character. Distribution type. Either ‘"LP3"’ (default) or
          ‘"GEV"’.

strat_dist: Character. Probability space for stratification bins.
          Passed to ‘stratified_sampler()’. One of ‘"ev1"’ (default),
          ‘"normal"’, or ‘"uniform"’. See ‘stratified_sampler()’ for
          details.

    Nbin: Integer. Number of stratified bins. Default is ‘50’.

  Mevent: Integer. Number of events per bin. Default is ‘200’.

Value:

     A list containing:

     flow Matrix of sampled flow values [Mevent x Nbin]

     nbins Number of stratified bins

     mevents Number of events per bin

     weights Probability weights for each bin from
          ‘stratified_sampler()’

See Also:

     ‘stratified_sampler()’, ‘qp3()’, ‘rfa_simulate()’

Examples:

     # Single parameter set (posterior mode)
     params <- c(4.85, 0.39, -0.15)
     result <- flow_frequency_sampler(params, freq_dist = "LP3",
                                      Nbin = 20, Mevent = 500)
     # Dimensions of result
     dim(result$flow)  # 500 x 20
     
     # Distribution of Sampled Flows
     hist(result$flow)
     
\end{verbatim}

\subsection{\texorpdfstring{\texttt{flow\_frequency\_sampler\_expected()}}{flow\_frequency\_sampler\_expected()}}\label{flow_frequency_sampler_expected}

\begin{verbatim}
Flow Frequency Sampler (Expected Only)

Description:

     Generates a stratified matrix of flow values by pairing each
     z-ordinate with a different parameter set from the posterior
     distribution. This collapses the nested Monte Carlo structure into
     a single pass, simultaneously sampling natural variability (via
     stratified z-ordinates) and knowledge uncertainty (via varying
     parameters). Used internally by ‘rfa_simulate()’ for expected-only
     mode.

Usage:

     flow_frequency_sampler_expected(
       bestfit_params,
       freq_dist = "LP3",
       strat_dist = "ev1",
       Nbin = NULL,
       Mevent = NULL
     )
     
Arguments:

bestfit_params: Data frame of distribution parameters from RMC-BestFit.
          Must have ‘Nbin * Mevent’ rows (one parameter set per
          z-ordinate). For LP3: columns are mean (log), sd (log), skew
          (log). For GEV: columns are location, scale, shape.

freq_dist: Character. Distribution type. Either ‘"LP3"’ (default) or
          ‘"GEV"’.

strat_dist: Character. Probability space for stratification bins.
          Passed to ‘stratified_sampler()’. One of ‘"ev1"’ (default),
          ‘"normal"’, or ‘"uniform"’. See ‘stratified_sampler()’ for
          details.

    Nbin: Integer. Number of stratified bins. Default is ‘50’.

  Mevent: Integer. Number of events per bin. Default is ‘200’.

Value:

     A list containing:

     flow Matrix of sampled flow values [Mevent x Nbin]

     nbins Number of stratified bins

     mevents Number of events per bin

     weights Probability weights for each bin from
          ‘stratified_sampler()’

See Also:

     ‘stratified_sampler()’, ‘qp3()’, ‘flow_frequency_sampler()’,
     ‘rfa_simulate()’

Examples:

     # Using a pre-loaded parameter set (all 10,000 parameter sets)
     result <- flow_frequency_sampler_expected(jmd_bf_parameter_sets,
                                               Nbin = 20, Mevent = 500)
     dim(result$flow)  # 500 x 20
     
     # Using bootstrapped parameter samples
     jmd_samples <- bootstrap_vfc(
       c(jmd_vfc_parameters$mean_log,
         jmd_vfc_parameters$sd_log,
         jmd_vfc_parameters$skew_log),
       dist = "LP3",
       ERL  = jmd_vfc_parameters$erl)
     
     jmd_result <- flow_frequency_sampler_expected(
       jmd_samples$params,
       freq_dist = "LP3",
       Nbin      = 20,
       Mevent    = 500)
     
\end{verbatim}

\subsection{\texorpdfstring{\texttt{hydrograph\_setup()}}{hydrograph\_setup()}}\label{hydrograph_setup}

\begin{verbatim}
Hydrograph Setup for RFA Simulation

Description:

     Prepares hydrograph data frames copied from RMC-RFA for use in RFA
     simulation. Converts date/time columns and adds sequential hour
     and hydrograph ID columns.

Usage:

     hydrograph_setup(
       ...,
       critical_duration = NULL,
       routing_days = NULL,
       weights = NULL
     )
     
Arguments:

     ...: Data frame representing an input hydrograph with columns:
          Ord, Date, Time, Flow (cfs). Copied directly from RMC-RFA.
          Date & Time should be ‘class() = "character"’

critical_duration: Critical duration in days.

routing_days: Desired length of routing simulation in days.

 weights: Optional numeric vector of sampling weights for each
          hydrograph. Must be the same length as the number of input
          hydrographs. Weights are normalized to probabilities
          internally. If NULL (default), all hydrographs are weighted
          equally.

Value:

     A list of formatted hydrograph data frames, each containing:

     datetime Date-time (POSIXct)

     hour Hours from start of event

     inflow Inflow (cfs)

     hydrograph_num Hydrograph ID number for sampling

     obs_vol Max n-day inflow volume (stored as an attribute of each
          dataframe)

     dt Hydrograph timestep (delta time, dt) in hours (stored as an
          attribute of each dataframe)

     The returned list also has a ‘probs’ attribute containing the
     normalized sampling probabilities derived from ‘weights’.

Examples:

     # Setup with equal weights (default)
     hydros <- hydrograph_setup(jmd_hydro_apr1999,
                                jmd_hydro_jun1965,
                                jmd_hydro_pmf,
                                critical_duration = 2,
                                routing_days = 10)
     
     # Setup with custom weights (PMF 3x more likely to be sampled)
     hydros <- hydrograph_setup(jmd_hydro_apr1999,
                                jmd_hydro_jun1965,
                                jmd_hydro_pmf,
                                critical_duration = 2,
                                routing_days = 10,
                                weights = c(1, 1, 3))
     
     # view normalized probabilities
     attr(hydros, "probs")
     
\end{verbatim}

\subsection{\texorpdfstring{\texttt{scale\_hydrograph()}}{scale\_hydrograph()}}\label{scale_hydrograph}

\begin{verbatim}
Scale Hydrograph

Description:

     Scales inflow hydrograph given a sampled inflow volume from the
     volume-frequency curve. Scale factor is defined by the sampled
     inflow volume and the maximum volume of the corresponding
     duration. If the native hydrograph timestep differs from the
     target routing timestep, resampling is applied prior to scaling:
     finer-to-coarser uses block averaging; coarser-to-finer uses
     linear interpolation.

Usage:

     scale_hydrograph(
       hydrograph_shape,
       observed_volume,
       sampled_volume,
       routing_dt = 1
     )
     
Arguments:

hydrograph_shape: Data frame with two columns: time (hours) and inflow
          (cfs). Must be in that order.

observed_volume: Maximum n-day inflow volume from input hydrograph
          shape.

sampled_volume: Sampled inflow volume from volume-frequency curve.

routing_dt: Target routing timestep in hours. Supported values are
          ‘0.25’ (15-min), ‘1’ (1-hour, default), ‘6’ (6-hour), and
          ‘24’ (24-hour).

Value:

     A data frame with two columns: ‘time_hrs’ and ‘inflow_cfs’, at the
     target routing timestep, with time starting at hour 0.

Examples:

     # Example hydrograph. Requires pre-processing
     hydro_example <- hydrograph_setup(jmd_hydro_jun1965_15min, critical_duration = 2, routing_days = 10)
     hydrograph_shape <- hydro_example[[1]][, 2:3]
     
     # Default 1-hour routing timestep
     scaled <- scale_hydrograph(hydrograph_shape,
                                observed_volume = 50000,
                                sampled_volume = 55000)
     
     # 15-min input hydrograph, 1-hour routing timestep
     scaled <- scale_hydrograph(hydrograph_shape,
                                observed_volume = 50000,
                                sampled_volume = 55000,
                                routing_dt = 1)
     
\end{verbatim}

\subsection{\texorpdfstring{\texttt{mod\_puls\_routing()}}{mod\_puls\_routing()}}\label{mod_puls_routing}

\begin{verbatim}
Modified Puls Reservoir Routing

Description:

     Performs Modified Puls (level pool) routing of inflow hydrograph
     given a defined reservoir geometry (Stage (ft), Storage (ac-ft),
     Discharge (cfs)).

Usage:

     mod_puls_routing(resmodel_df, inflow_df, initial_elev, full_results = FALSE)
     
Arguments:

resmodel_df: Data frame with three columns: elevation/stage (ft),
          storage (acre-feet), and discharge (cfs). Must be in that
          order.

inflow_df: Data frame with two columns: time (hours) and inflow (cfs).
          Must be in that order.

initial_elev: Starting water surface elevation in feet.

full_results: Logical. If ‘FALSE’ (default), returns only peak stage
          and discharge. If ‘TRUE’, returns the complete routing
          result.

Value:

     If ‘full_results = FALSE’, a named numeric vector with
     ‘peak_stage_ft’ and ‘peak_discharge_cfs’. If ‘full_results =
     TRUE’, a data frame with columns: ‘time_hr’, ‘inflow_cfs’,
     ‘elevation_ft’, ‘storage_acft’, and ‘outflow_cfs’.

References:

     Chow, V.T. (1959). Open-Channel Hydraulics. McGraw-Hill.

Examples:

     # Example hydrograph. Requires pre-processing
     hydro_example <- hydrograph_setup(jmd_hydro_jun1965_15min,
                                        critical_duration = 2,
                                        routing_days = 10)
     hydrograph_shape <- hydro_example[[1]][, 2:3]
     scaled_hydrograph <- scale_hydrograph(hydrograph_shape,
                                           observed_volume = 50000,
                                           sampled_volume = 55000)
     
     # Peak values only
     mod_puls_routing(jmd_resmodel, scaled_hydrograph, initial_elev = 3830)
     
     # Full routing table
     jmd_full_routing <- mod_puls_routing(jmd_resmodel,
                                          scaled_hydrograph,
                                          initial_elev = 3830,
                                          full_results = TRUE)
     head(jmd_full_routing)
     
\end{verbatim}

\subsection{\texorpdfstring{\texttt{stage\_frequency\_curve()}}{stage\_frequency\_curve()}}\label{stage_frequency_curve}

\begin{verbatim}
Compute Stage-Frequency Curve

Description:

     Converts a matrix of routed peak stages into a stage-frequency
     curve using the law of total probability across stratified bins.
     For each candidate stage, the conditional exceedance probability
     within each bin is weighted by the bin probability and summed to
     produce the unconditional annual exceedance probability (AEP).

Usage:

     stage_frequency_curve(peakStage, weights, stage_bins = 1000)
     
Arguments:

peakStage: Matrix of peak stages [Mevents x Nbins] from routing.

 weights: Numeric vector of bin weights from ‘stratified_sampler()’.
          Must have length equal to ‘ncol(peakStage)’.

stage_bins: Numeric value of exceedance stages used to compute
          stage-frequency in each bin. Default is 1000.

Value:

     A data frame with columns:

     stage Evaluated stage values

     AEP Annual exceedance probability at each stage

See Also:

     ‘stratified_sampler()’, ‘flow_frequency_sampler()’,
     ‘rfa_simulate()’
\end{verbatim}

\section{Rejection Sampling}\label{rejection-sampling}

\subsection{\texorpdfstring{\texttt{pmf\_stage\_lognormal()}}{pmf\_stage\_lognormal()}}\label{pmf_stage_lognormal}

\begin{verbatim}
Parameterize a Three-Parameter Lognormal PMF for Stage

Description:

     Computes the parameters of a three-parameter lognormal
     distribution for use as a probabilistic maximum stage (PMF) in
     rejection sampling. The shift parameter defines the lower or upper
     bound of the distribution (typically, the lower value). The
     function supports two modes: (1) supplying a best estimate and
     assuming sigma, or (2) supplying a best estimate, low, and high to
     solve for sigma numerically.

Usage:

     pmf_stage_lognormal(
       pmf_shift,
       pmf_mean,
       pmf_sigma = NULL,
       pmf_low = NULL,
       pmf_high = NULL
     )
     
Arguments:

pmf_shift: Numeric. Assumed lower-bound of the lognormal distribution.
          This will define the shift value. Must be less than all of
          ‘pmf_mean’, ‘pmf_low’, and ‘pmf_high’.

pmf_mean: Numeric. Assumed mean of the shifted PMF stage distribution.

pmf_sigma: Numeric. Assumed standard deviation on the log scale.
          Required if ‘pmf_low’ and ‘pmf_high’ are not supplied.
          Defaults to ‘NULL’. Suggested value = 0.5.

 pmf_low: Numeric. Low estimate of PMF stage (assumed 5th percentile).
          Optional. If supplied, ‘pmf_high’ must also be supplied.

pmf_high: Numeric. High estimate of PMF stage (assumed 95th
          percentile). Optional. If supplied, ‘pmf_low’ must also be
          supplied.

Value:

     A named list with the following elements:

     pmf_shift Hard lower bound of the distribution.

     pmf_mean Assumed mean of shifted PMF stage distribution.

     pmf_sigma Standard deviation on the log scale.

     pmf_mu Derived location parameter on the log scale.

     pmf_p05 5th percentile of the PMF stage distribution.

     pmf_p95 95th percentile of the PMF stage distribution.

Examples:

     # Mode 1: assume sigma
     pmf <- pmf_stage_lognormal(pmf_shift = 239, pmf_mean = 241.9, pmf_sigma = 0.5)
     
     # Mode 2: solve for sigma from low/high estimates
     pmf <- pmf_stage_lognormal(pmf_shift = 239, pmf_mean = 241.9,
                                 pmf_low = 239, pmf_high = 245)
     
\end{verbatim}

\subsection{\texorpdfstring{\texttt{rejection\_sampling\_stage()}}{rejection\_sampling\_stage()}}\label{rejection_sampling_stage}

\begin{verbatim}
Rejection Sampling of Stage Bounded by a Probabilistic Maximum Stage

Description:

     Draws ‘n_samples’ stage values from a stage-frequency curve,
     rejecting any draws that exceed a probabilistic maximum stage
     (PMF) sampled from a three-parameter lognormal distribution.
     Accepted samples are ranked and assigned Weibull plotting
     positions for use in frequency analysis.

Usage:

     rejection_sampling_stage(
       pmf_stage_LN,
       stage_freq_df = NULL,
       aep = NULL,
       stage = NULL,
       n_samples = 1e+07
     )
     
Arguments:

pmf_stage_LN: A named list returned by ‘pmf_stage_lognormal’ containing
          the lognormal PMF parameters.

stage_freq_df: A data frame with AEP in column 1 and stage in column 2.
          Mutually exclusive with ‘aep’ and ‘stage’.

     aep: Numeric vector of annual exceedance probabilities. Must be
          supplied with ‘stage’. Mutually exclusive with
          ‘stage_freq_df’.

   stage: Numeric vector of stages corresponding to ‘aep’. Must be
          supplied with ‘aep’. Mutually exclusive with ‘stage_freq_df’.

n_samples: Integer. Number of samples to draw. Defaults to ‘1e7’.

Value:

     A data frame of thinned accepted samples with columns for z-score
     and stage, suitable for plotting a probabilistically bounded
     stage-frequency curve. Output is produced by ‘thin_samples()’.

See Also:

     ‘pmf_stage_lognormal’

Examples:

     ## Not run:
     
     pmf_ln <- pmf_stage_lognormal(pmf_shift = 239, pmf_mean = 241.9, pmf_sigma = 0.5)
     
     result <- rejection_sampling_stage(
       pmf_stage_LN  = pmf_ln,
       stage_freq_df = my_stage_freq_df,
       n_samples     = 1e7
     )
     ## End(Not run)
     
\end{verbatim}

\section{Additional Utilities}\label{additional-utilities}

\subsection{\texorpdfstring{\texttt{bootstrap\_vfc()}}{bootstrap\_vfc()}}\label{bootstrap_vfc}

\begin{verbatim}
Bootstrap option for VFC (RFA-style w/ ERL)

Description:

     Generates matrix of parameters using posterior mode/mean and ERL

Usage:

     bootstrap_vfc(bestfit_postmode, dist = "LP3", ERL = 150, Nboots = 10000)
     
Arguments:

bestfit_postmode: Vector of distribution parameters from RMC-BestFit.
          For LP3: columns are mean (log), sd (log), skew (log). For
          GEV: columns are location, scale, shape.

    dist: Distribution type. Either ‘"LP3"’ (default) or ‘"GEV"’.

     ERL: psuedo effective record length for bootstrapping

  Nboots: number of bootstraps

Value:

     A list containing:

     params Matrix of bootstrapped parameters

     dist Distribution selected (LP3 or GEV)

     postmode Posterior mode parameters provided to bootstrap

     ERL pseudo effective record length provided to bootstrap

See Also:

     ‘flow_frequency_sampler()’

Examples:

     # Sample using JMD VFC parameters
     jmd_samples <- bootstrap_vfc(c(jmd_vfc_parameters$mean_log,
                                     jmd_vfc_parameters$sd_log,
                                     jmd_vfc_parameters$skew_log),
                                     dist = "LP3",
                                     ERL = jmd_vfc_parameters$erl)
     
     # GEV Example with parent distribution parameters
     gev_example <- c(3.0, 1.0, -0.1)
     gev_samples <- bootstrap_vfc(gev_example,
                                   dist = "GEV",
                                   ERL = 200,
                                   Nboots = 5000)
     hist(gev_samples$params[,1])
     
     # LP3 Example with parent distribution parameters
     lp3_example <- c(3.5, 0.22, 0.1)
     lp3_samples <- bootstrap_vfc(lp3_example,
                                   dist = "LP3",
                                   ERL = 300,
                                   Nboots = 1000)
     hist(lp3_samples$params[,3])
     
\end{verbatim}

\subsection{\texorpdfstring{\texttt{qp3()}}{qp3()}}\label{qp3}

\begin{verbatim}
Pearson Type III Inverse CDF (Quantile Function)

Description:

     Computes quantiles from a Pearson Type III distribution given
     probabilities and distribution parameters (mean, standard
     deviation, skewness).

Usage:

     qp3(p, mu, sigma, gamma)
     
Arguments:

       p: Vector of probabilities (between 0 and 1).

      mu: Mean of the distribution.

   sigma: Standard deviation of the distribution.

   gamma: Skewness coefficient of the distribution.

Details:

     NOTE: This is the non-vectorized version of the function (it will
     be used in a loop) Vectorized version is part of future
     development.

     The Pearson Type III distribution is parameterized by mean (‘mu’),
     standard deviation (‘sigma’), and skewness (‘gamma’). These are
     converted internally to location, scale, and shape parameters.

     When skewness is near zero (‘abs(gamma) < 1E-3’), the normal
     distribution is used as an approximation.

Value:

     Vector of quantiles corresponding to the input probabilities.

Examples:

     # Single quantile
     qp3(0.99, mu = 10, sigma = 2, gamma = 0.5)
     
     # Multiple quantiles
     qp3(c(0.5, 0.9, 0.99), mu = 10, sigma = 2, gamma = 0.5)
     
     # Use with JMD volume-frequency parameters
     qp3(0.99,
         mu = jmd_vfc_parameters$mean_log,
         sigma = jmd_vfc_parameters$sd_log,
         gamma = jmd_vfc_parameters$skew_log)
     
\end{verbatim}

\subsection{\texorpdfstring{\texttt{aep2stage()}}{aep2stage()}}\label{aep2stage}

\begin{verbatim}
Interpolate Stage from AEP

Description:

     Given a stage-frequency curve, interpolates stage values at target
     AEPs. Interpolation is performed in standard normal (z-variate)
     space for better behavior across orders of magnitude in AEP.

Usage:

     aep2stage(aep, stage, interp_aep)
     
Arguments:

     aep: Numeric vector of annual exceedance probabilities from the
          known curve.

   stage: Numeric vector of stage values corresponding to ‘aep’.

interp_aep: Numeric vector of target AEPs to interpolate at.

Value:

     Numeric vector of interpolated stage values (same length as
     ‘interp_aep’).

See Also:

     ‘stage2aep’ for the inverse operation.

Examples:

     # Get stage at the 1.12% and 0.175% AEP
     aep2stage(jmd_rfa_expected$AEP, jmd_rfa_expected$Expected, c(0.0112,1.75E-3))
     
\end{verbatim}

\subsection{\texorpdfstring{\texttt{stage2aep()}}{stage2aep()}}\label{stage2aep}

\begin{verbatim}
Interpolate AEP from Stage

Description:

     Given a stage-frequency curve, interpolates AEP values at target
     stages. Interpolation is performed in standard normal (z-variate)
     space for better behavior across orders of magnitude in AEP.

Usage:

     stage2aep(aep, stage, interp_stage)
     
Arguments:

     aep: Numeric vector of annual exceedance probabilities from the
          known curve.

   stage: Numeric vector of stage values corresponding to ‘aep’.

interp_stage: Numeric vector of target stage values to interpolate at.

Details:

     Stage-frequency curves typically flatten at very rare AEPs once
     the reservoir reaches the reservoir model discharge capacity. This
     produces tied stage values. ‘approx()’ handles these by averaging
     the corresponding AEPs, which may emit a "collapsing to unique 'x'
     values" warning. This is expected behavior.

Value:

     Numeric vector of interpolated AEP values (same length as
     ‘interp_stage’).

See Also:

     ‘aep2stage’ for the inverse operation.

Examples:

     # Get AEP at stage 3862.0-ft and 3872-ft
     stage2aep(jmd_rfa_expected$AEP, jmd_rfa_expected$Expected, c(3862.0, 3872))
     
\end{verbatim}

\subsection{\texorpdfstring{\texttt{interpolate\_aep\_matrix()}}{interpolate\_aep\_matrix()}}\label{interpolate_aep_matrix}

\begin{verbatim}
Interpolate AEPs at Target Stages Across Realizations

Description:

     For each realization's stage-frequency curve, interpolates the AEP
     at every target stage value. Interpolation is performed in
     standard normal (z-variate) space.

Usage:

     interpolate_aep_matrix(results_list, target_stages)
     
Arguments:

results_list: List of realization data frames, each with columns
          ‘stage’ and ‘AEP’.

target_stages: Numeric vector of stage values to interpolate at.

Value:

     Matrix of AEP values with rows = target stages and columns =
     realizations.
\end{verbatim}

\subsection{\texorpdfstring{\texttt{interpolate\_stage\_matrix()}}{interpolate\_stage\_matrix()}}\label{interpolate_stage_matrix}

\begin{verbatim}
Interpolate Stages at Target AEPs Across Realizations

Description:

     For each realization's stage-frequency curve, interpolates the
     stage at every target AEP value. Interpolation is performed in
     standard normal (z-variate) space.

Usage:

     interpolate_stage_matrix(results_list, target_aeps)
     
Arguments:

results_list: List of realization data frames, each with columns
          ‘stage’ and ‘AEP’.

target_aeps: Numeric vector of AEP values to interpolate at.

Value:

     Matrix of stage values with rows = target AEPs and columns =
     realizations.
\end{verbatim}

\subsection{\texorpdfstring{\texttt{power\_function()}}{power\_function()}}\label{power_function}

\begin{verbatim}
Power Function

Description:

     Raises 10 to the 1st through the nth power

Usage:

     power_function(x)
     
Arguments:

       x: Highest power of 10 to use

Value:

     Vector of 10 raised sequentially from the 1st to the nth power.
\end{verbatim}

\subsection{\texorpdfstring{\texttt{theme\_rfar\_conceptual()}}{theme\_rfar\_conceptual()}}\label{theme_rfar_conceptual}

\begin{verbatim}
rfaR ggplot2 Theme for Conceptual Plots

Description:

     A custom ‘ggplot2’ theme used in the rfaR conceptual realization
     vignette. Based on ‘theme_bw’ with tightened text sizes, italic
     titles, and the legend suppressed. Intended for figures that
     emphasize curves over annotation density.

Usage:

     theme_rfar_conceptual()
     
Value:

     A ‘ggplot2’ theme object that can be added to a ‘ggplot’ via ‘+’.

See Also:

     ‘theme_bw’

Examples:

     library(ggplot2)
     ggplot(jmd_rfa_expected, aes(x = AEP, y = Expected)) +
       geom_line() +
       scale_x_continuous(transform = c("log10", "reverse")) +
       labs(title = "JMD Expected Stage-Frequency",
            subtitle = "Demonstrating theme_rfar_conceptual()",
            x = "AEP", y = "Stage (ft)") +
       theme_rfar_conceptual()
     
\end{verbatim}

\protect\phantomsection\label{refs}
\begin{CSLReferences}{1}{1}
\bibitem[\citeproctext]{ref-aven_society_nodate}
Aven, Terje, Yakov Ben-Haim, Henning Boje Andersen, et al. n.d.
\emph{Society for {Risk} {Analysis} {Glossary}}.

\bibitem[\citeproctext]{ref-noaa_nws_prsk2_2026}
National Oceanic and Atmospheric Administration. 2026. \emph{National
Water Prediction Service: Official Hydrograph for Gauge PRSK2}. National
Weather Service. \url{https://water.noaa.gov/gauges/prsk2\#hg=official}.

\end{CSLReferences}




\end{document}
